% §2.2 The Asymmetric Binary Channel
%
% Cross-reference resolution status:
%
% RESOLVED:
%   \Cref{sec:three-faces}        -- chapter 1 §1.1
%   \Cref{sec:bsc}, \Cref{def:bsc}-- chapter 2 §2.1
%   \Cref{sec:model-order}        -- chapter 2 §2.4
%   \Cref{sec:channel-composition}-- chapter 2 §2.3
%
% Monograph register: defines the general asymmetric channel and the
% Z-channel (with the Bloom-filter connection), gives one instance, and
% reconciles model order consistently with §2.4 (BSC is order 1; the
% asymmetric channel is the general order-2 case). The eps=omega
% measure-zero point is salvaged as a remark.

\section{The Asymmetric Binary Channel}
\label{sec:asymmetric-channel}

The BSC's equal error rates are rarely satisfied in practice: a spam filter, a
Bloom filter, and a medical screening test all treat their two error directions
differently. Allowing the false-positive rate\index{false-positive rate} $\fprate$ and the false-negative
rate\index{false-negative rate} $\fnrate$ to vary independently gives the general binary channel,
the atom of two-sided error\index{two-sided error}.

\begin{definition}[General asymmetric binary channel]
\label{def:asymmetric-channel}
An \emph{asymmetric binary channel} with rates $\fprate, \fnrate \in [0,1)$ is
a random map $\phi : \{\top, \bot\} \to \{\top, \bot\}$ with
$\Pr\{\phi(\top) \neq \top\} = \fnrate$ and $\Pr\{\phi(\bot) \neq \bot\} =
\fprate$, and channel matrix
\[
  Q \;=\;
  \begin{pmatrix}
    1 - \fnrate & \fnrate \\
    \fprate     & 1 - \fprate
  \end{pmatrix}
\]
in the row convention of \Cref{sec:bsc}.
\end{definition}

Setting $\fnrate = \fprate$ recovers the BSC of \Cref{def:bsc}. Setting
$\fnrate = 0$ makes row 1 equal to $(1, 0)$, so a positive input is always
reported correctly: this is the Z-channel. \Cref{fig:binary-channel} draws the
three cases. The asymmetric channel\index{channel!asymmetric binary} is the
standard binary channel of information theory. Read as a model of how a
recorded indicator departs from the true one, it is also the misclassification
model of $2 \times 2$ tables and, when the noise is introduced deliberately,
Warner's randomized response\index{randomized response} (see the
\hyperref[sec:bib-notes-ch2]{Notes and Lineage} at the end of this chapter).

\begin{figure}[tbp]
\centering
\begin{tikzpicture}[x=24mm,y=13mm,
  sym/.style={circle,draw,inner sep=0.5pt,minimum size=6mm},
  lab/.style={font=\footnotesize,inner sep=1.5pt},
  arr/.style={-{Stealth[length=2mm]}}]
  % (a) BSC
  \begin{scope}[shift={(0,0)}]
    \node[sym] (it) at (0,1) {$\top$};
    \node[sym] (ib) at (0,0) {$\bot$};
    \node[sym] (ot) at (1,1) {$\top$};
    \node[sym] (ob) at (1,0) {$\bot$};
    \draw[arr] (it) -- node[lab,above] {$1-\fprate$} (ot);
    \draw[arr] (ib) -- node[lab,below] {$1-\fprate$} (ob);
    \draw[arr] (it) -- node[lab,pos=0.78,anchor=south west] {$\fprate$} (ob);
    \draw[arr] (ib) -- node[lab,pos=0.78,anchor=north west] {$\fprate$} (ot);
    \node[font=\footnotesize] at (0.5,-0.75) {(a) binary symmetric};
  \end{scope}
  % (b) asymmetric
  \begin{scope}[shift={(1.7,0)}]
    \node[sym] (it) at (0,1) {$\top$};
    \node[sym] (ib) at (0,0) {$\bot$};
    \node[sym] (ot) at (1,1) {$\top$};
    \node[sym] (ob) at (1,0) {$\bot$};
    \draw[arr] (it) -- node[lab,above] {$1-\fnrate$} (ot);
    \draw[arr] (ib) -- node[lab,below] {$1-\fprate$} (ob);
    \draw[arr] (it) -- node[lab,pos=0.78,anchor=south west] {$\fnrate$} (ob);
    \draw[arr] (ib) -- node[lab,pos=0.78,anchor=north west] {$\fprate$} (ot);
    \node[font=\footnotesize] at (0.5,-0.75) {(b) asymmetric};
  \end{scope}
  % (c) Z-channel
  \begin{scope}[shift={(3.4,0)}]
    \node[sym] (it) at (0,1) {$\top$};
    \node[sym] (ib) at (0,0) {$\bot$};
    \node[sym] (ot) at (1,1) {$\top$};
    \node[sym] (ob) at (1,0) {$\bot$};
    \draw[arr] (it) -- node[lab,above] {$1$} (ot);
    \draw[arr] (ib) -- node[lab,below] {$1-\fprate$} (ob);
    \draw[arr] (ib) -- node[lab,pos=0.78,anchor=north west] {$\fprate$} (ot);
    \node[font=\footnotesize] at (0.5,-0.75) {(c) Z-channel};
  \end{scope}
  \node[font=\footnotesize] at (0,1.55) {truth};
  \node[font=\footnotesize] at (1,1.55) {report};
\end{tikzpicture}
\caption[The binary channels: BSC, asymmetric channel, Z-channel]{The binary channels of Part~I. Each arrow carries the probability of
the report given the truth; the rows of the channel matrix $Q$ are the arrow
weights leaving one input. (a) The BSC has one crossover $\fprate$ in both
directions (model order 1). (b) The asymmetric channel has a false-positive
rate $\fprate$ and a false-negative rate $\fnrate$ (model order 2). (c) The
Z-channel fixes $\fnrate = 0$: a report of $\bot$ is always correct, the
one-sided profile of a Bloom filter query.}
\label{fig:binary-channel}
\end{figure}

\begin{definition}[Z-channel]
\label{def:z-channel}
A \emph{Z-channel}\index{channel!Z-channel} with false-positive rate $\fprate \in (0,1)$ is an
asymmetric binary channel with $\fnrate = 0$, channel matrix
\[
  Q_Z \;=\;
  \begin{pmatrix}
    1       & 0 \\
    \fprate & 1 - \fprate
  \end{pmatrix}.
\]
\end{definition}

The Z-channel makes false negatives impossible while admitting false positives
at rate $\fprate$ (one-sided error\index{one-sided error}), so an output of $\bot$ is fully reliable and an output of
$\top$ is uncertain.\footnote{Information-theory texts often take the opposite
asymmetry ($\fnrate > 0$, $\fprate = 0$); the two are equivalent up to
relabeling. The convention here aligns with positive approximate sets, where
the false-positive direction is the natural error mode.} This reliable-no,
uncertain-yes profile is exactly the Bloom filter: it answers set membership
with no false negatives ($x \in S$ always returns $\top$) and a tunable
false-positive rate ($x \notin S$ may still return $\top$), so a Bloom
filter\index{Bloom filter!as a channel} query is a Z-channel, $\fnrate = 0$
with $\fprate$ set by its bit-width and hash count (conditionally on the
realized filter; \Cref{ch:approximate-sets} treats the randomness of the
rate itself). The doubly-biased coin, the screening test (small $\fnrate$, larger
$\fprate$) and its confirmatory counterpart (the reverse), and any system tuned
by a decision threshold are all asymmetric channels; as the threshold sweeps,
$(\fprate, \fnrate)$ traces a curve in $[0,1]^2$, each point a distinct channel,
which is the ROC curve re-axed.

For one instance, take $\fprate = 0.05$, $\fnrate = 0.10$:
\[
  Q \;=\;
  \begin{pmatrix}
    0.90 & 0.10 \\
    0.05 & 0.95
  \end{pmatrix},
\]
with false negatives twice as likely as false positives. Under a uniform input
prior its accuracy is the diagonal average $\tfrac{1}{2}(0.90 + 0.95) = 0.925$.
A BSC at $\fprate = \fnrate = 0.075$ achieves the same accuracy with both
off-diagonal entries equal, so accuracy alone does not see the asymmetry; the
distinction matters once the prior is non-uniform or the two errors carry
different costs, which is the subject of \Cref{ch:classification-measures}.

\subsection*{The general order-2 case}

The asymmetric channel is the general model-order-2\index{model order} binary channel: the input
domain partitions into $B_1 = \{\top\}$ and $B_2 = \{\bot\}$, each block
carrying its own rate, $\fnrate$ for $B_1$ and $\fprate$ for $B_2$. The BSC
imposes $\fnrate = \fprate$, which merges the two blocks into one and drops the
model order to 1 (\Cref{sec:model-order}); the Z-channel keeps the blocks
distinct ($\fnrate = 0 \neq \fprate$) and so remains order 2 with a single free
parameter, the first sign that model order and parameter count are separate
quantities\index{model order!versus parameter count}. \Cref{sec:model-order} makes the distinction precise and proves the
Boolean ceiling: with only two input elements, no binary channel exceeds model
order 2.

\begin{remark}[The BSC is a measure-zero special case]
\label{rem:bsc-measure-zero}
In the parameter square $[0,1)^2$ of pairs $(\fprate, \fnrate)$, the binary
symmetric channels are exactly the diagonal $\fprate = \fnrate$, a
one-dimensional, measure-zero subset. A generic asymmetric channel is not a
BSC, and an arbitrarily small perturbation off the diagonal turns a BSC into a
two-block, model-order-2 channel. Symmetry is the exception in the space of
binary channels, not the rule, which is why the book works in the order-2 frame
and treats the BSC as the constrained case rather than the reverse.
\end{remark}
